% !PRINTED 171
\weberpage{171}

%<<CONT
die Wurzeln der Gleichung (4), so ergeben sich aus (5) die entsprechenden Werthe von $y$:
\weq{7}{y_1 = \varphi(x_1), \quad y_2 = \varphi(x_2) \ldots y_n = \varphi(x_n),}
und das Product
\weq{8}{\Phi(y) = (y - y_1)(y - y_2) \ldots (y - y_n)}
% [corrected 2026-08-07, 8x zoom strip 2/8] Three lines letterspaced. "ist eine" is set normally
% and the run begins at "symmetrische", carrying through the reference "(6)" and the closing
% full stop. Claimed by the translation pass and independently by the sperrung pass.
ist eine \gesperrt{symmetrische Function der Grössen (6), die sich also rational durch die Co\"efficienten von $f(x)$ ausdrücken lässt.}

In dieser Weise dargestellt, möge $\Phi(y)$ den Ausdruck haben
\weq{9}{\Phi(y) = y^n + b_1 y^{n-1} + b_2 y^{n-2} + \cdots + b_{n-1} y + b_n.}

Die Co\"efficienten $b_\nu$ werden nach \S.\,7 durch die $y_i$ bestimmt mittelst der Formeln
\[
b_1 = -\Sigma y_1, \quad b_2 = \Sigma y_1 y_2, \quad b_3 = -\Sigma y_1 y_2 y_3 \ldots,
\]
\gesperrt{woraus sich ergiebt, dass $b_\nu$ eine ganze rationale und homogene Function $\nu^{\text{ten}}$ Grades von den Variablen $\alpha_0, \alpha_1 \ldots \alpha_{n-1}$ ist.}

$\Phi(y)$ ist nichts Anderes als die Resultante der beiden Gleichungen
\weq{10}{f(x) = 0, \quad \varphi(x) - y = 0,}
und nach \S.\,49 kann man auch umgekehrt die gemeinsame Wurzel dieser beiden Gleichungen (10), wenn $y$ einem der Werthe $y_1, y_2 \ldots y_n$ gleich wird, rational durch die Co\"efficienten von (10), d.\,h. rational durch $a_i$, $\alpha_i$, $y$ ausdrücken, vorausgesetzt, dass \gesperrt{nur eine solche gemeinsame Wurzel vorhanden ist.}

Dies wird dann eintreten, wenn die $n$ Werthe $y_1, y_2 \ldots y_n$ von einander verschieden sind. Dann erhält man einen Ausdruck
\[
x = \beta_0 + \beta_1 y + \beta_2 y^2 + \cdots + \beta_{n-1} y^{n-1},
\]
der gültig ist, wenn für $x$ einer der Werthe (6) und für $y$ der zugehörige Werth (7) gesetzt wird. Die Grössen $\beta_i$ sind aus den $a_i$ und $\alpha_i$ rational zusammengesetzt.

Kommen aber unter den Werthen $y_1, y_2 \ldots y_n$ mehrere gleiche vor, so erfordert die Bestimmung der zugehörigen Werthe der $x$ noch die Auflösung einer Gleichung des entsprechenden Grades. \gesperrt{Dieser Grad ist aber immer kleiner als $n$}, da die Function $\varphi(x)$ höchstens für $n-1$ Werthe von $x$ denselben Werth $y$ erhalten kann. Wenn also die Gleichung
% [source-original mathematical caveat, confirmed 2026-08-25 against scan p. 197 / printed p. 171]
% The assertion omits the constant-phi case: if phi is constant, all n roots x_i can have the
% same y-value. The printed claim is preserved rather than silently qualified.
\weq{11}{\Phi(y) = 0}
%CONT>>
